% TOP_PROBLEM: {"problemId": "problem.uniqueness-of-the-kpz-fixed-point-among-scaling-covariant-local-fields", "canonicalTitle": "Uniqueness of the KPZ Fixed Point Among Scaling-Covariant Local Fields", "releaseRank": 452, "primaryDomain": "probability_ergodic_dynamics", "primaryDomainLabel": "Probability, ergodic theory & dynamics", "displayStatus": "open", "releaseStatus": "open"}
% !TeX program = lualatex
% Definition and sources only; this file is not an LLM solution attempt.
\documentclass[11pt]{article}
\usepackage[margin=25mm]{geometry}
\usepackage{fontspec}
\setmainfont{Latin Modern Roman}
\usepackage{amsmath,amssymb,mathtools}
\usepackage{unicode-math}
\setmathfont{Latin Modern Math}
\usepackage{xurl,hyperref}
\hypersetup{colorlinks=true,urlcolor=blue}
\setcounter{secnumdepth}{0}
\setlength{\parindent}{0pt}
\setlength{\parskip}{5pt}
\setlength{\emergencystretch}{3em}
\begin{document}
\section*{\texorpdfstring{Uniqueness of the KPZ Fixed Point Among Scaling-Covariant Local Fields}{Uniqueness of the KPZ Fixed Point Among Scaling-Covariant Local Fields}}\label{452}
\subsection{Definitions and mathematical statement}
Use the state space $\mathrm{UC}$ of upper-semicontinuous profiles $h:{\mathbb{R}}\to{\mathbb{R}}\cup\{-\infty\}$, not identically $-\infty$, with $h(x)\le C(1+|x|)$ for some $C$. The KPZ fixed point $\mathfrak h(t,x;h_0)$ is the profile-valued process defined by Matetski--Quastel--Remenik's Definition 3.12 and transition formula (4.12).

The uniqueness question of Remark 4.12 specifies locality (4.19) and Theorem 4.5(i), (iii), (iv). Locality requires changes to $h_0$ outside neighborhoods of finitely many observation points to change their joint distribution by $o(t)$ as $t\downarrow0$. Scaling means, for $\alpha>0$,
\[
 \alpha\mathfrak h(\alpha^{-3}t,\alpha^{-2}x;
       \alpha^{-1}h_0(\alpha^2\cdot))
 \overset{\rm law}=\mathfrak h(t,x;h_0).
\]
The other two properties are skew time reversibility, interchanging initial and test profiles with sign reversal, and covariance under simultaneous spatial translation. Their full profile-event formulations are retained by the exact theorem reference. The target is uniqueness in law of the nontrivial field satisfying this package, not scaling invariance alone, and not uniqueness of solutions to the KPZ stochastic partial differential equation.
\par\smallskip\textit{Formulation and concept references:} \href{https://arxiv.org/html/1701.00018v3}{[S1]}.

\subsection{Short English statement}
Do the specified locality and symmetry requirements characterize the KPZ fixed point uniquely among nontrivial space-time fields?
\subsection{Sources}
\begin{itemize}
\item[S1]K. Matetski, J. Quastel, and D. Remenik, The KPZ fixed point, Remark 4.12, Acta Mathematica 2021.
\url{https://arxiv.org/html/1701.00018v3}.
\textit{Formulation source. Consulted 24 September 2026: Remark 4.12, locality (4.19), and Theorem 4.5(i), (iii), (iv).}
\end{itemize}
\end{document}
